\section{X25519 as an x-only byte-string function}

Curve25519 is a curve. X25519 is the RFC-defined function
\[
  \{0,1\}^{256}\times\{0,1\}^{256}
  \longrightarrow\{0,1\}^{256}
\]
that decodes a scalar and a Montgomery $u$-coordinate, evaluates x-only scalar
multiplication, and returns a 32-byte $u$-coordinate \cite{rfc7748}.

\subsection{Scalar and coordinate decoding}

The scalar is a 32-byte little-endian string. Before interpreting it as an
integer, X25519 applies the bit operation commonly called \emph{clamping}:

\begin{lstlisting}[language=C]
k[0]  &= 248;  /* clear bits 0, 1, and 2 */
k[31] &= 127;  /* clear bit 255 */
k[31] |= 64;   /* set bit 254 */
\end{lstlisting}

The resulting integer is exactly
\[
  k=2^{254}+8m,
  \qquad 0\le m\le2^{251}-1.
\]
It is divisible by $8$, its top processed bit is fixed, and it is not simply
the original bytes reduced modulo $\ell$. Clamping cannot create entropy.

An input $u$ is also decoded little-endian. X25519 masks the most significant
bit of the final byte and accepts noncanonical values, processing them as field
elements modulo $p$. Generic ``full point validation'' is not the X25519 input
rule: no $v$-coordinate is supplied, and the ladder intentionally handles both
the curve and its twist.

\subsection{The Montgomery ladder}

Conceptually, after processing a scalar prefix $r$, the ladder maintains
adjacent multiples
\[
  R_0=[r]P,
  \qquad
  R_1=[r+1]P.
\]
For the next bit, the abstract transition is
\[
\begin{array}{c|cc}
\text{bit}&R_0'&R_1'\\
\hline
0&[2]R_0&R_0+R_1\\
1&R_0+R_1&[2]R_1.
\end{array}
\]
One doubling and one differential addition occur for either bit. Constant-time
conditional swaps select the registers without a secret-dependent branch.

The implementation stores projective coordinates $X:Z$ representing
$u=X/Z$. This defers the costly field inversion until the end. Let $x_1$ be
the input $u$ and let $(X_2,Z_2)$ and $(X_3,Z_3)$ represent the two ladder
registers. One RFC ladder step uses
\begin{align*}
  A_0&=X_2+Z_2, & AA&=A_0^2,\\
  B_0&=X_2-Z_2, & BB&=B_0^2,\\
  E&=AA-BB,\\
  C_0&=X_3+Z_3, & D_0&=X_3-Z_3,\\
  DA&=D_0A_0, & CB&=C_0B_0,
\end{align*}
followed by
\begin{align*}
  X_3'&=(DA+CB)^2,&
  Z_3'&=x_1(DA-CB)^2,\\
  X_2'&=AA\,BB,&
  Z_2'&=E(AA+a_{24}E),
\end{align*}
where
\[
  a_{24}=\frac{486662-2}{4}=121665.
\]
All expressions are reduced modulo $p$. Another common, equivalent convention
uses $121666=(A+2)/4$ together with $BB$ in the last formula; mixing the
constant from one convention with the formula from the other is a bug.

After scanning bits $254$ through $0$, the output is
\[
  u=X_2Z_2^{p-2}\pmod p,
\]
encoded as 32 little-endian bytes. Fermat inversion also maps $Z_2=0$ to the
all-zero output in the RFC computation.

The fixed arithmetic pattern removes a common timing leak, but regularity is
not a proof against every cache, power, fault, or microarchitectural attack.
The companion Python follows the formulas for inspection and test vectors; its
branches and big integers are explicitly variable-time.

\subsection{X25519 Diffie--Hellman}

The base coordinate is encoded as the byte $9$ followed by 31 zero bytes.
Alice and Bob compute
\begin{align*}
  A&=\X(a,9), & B&=\X(b,9),\\
  S_A&=\X(a,B), & S_B&=\X(b,A).
\end{align*}
For honest inputs, $S_A=S_B$ because both encode $u([ab]G)$. RFC~7748
recommends deriving application keys from the shared output and both public
keys; an enclosing protocol must also specify how an all-zero output is
handled.
